
\subsection{Preparing for the Work}

Installment 5: “Folkebladet,” March 22, 1916

Besides what I pointed out in my previous article, another matter may also be mentioned in connection with intercessory prayer for a meeting: our prayer, and perhaps especially our intercession, is often so powerless, so little heaven-storming and heartfelt.

So often there is lacking that intensity and persistence with which we pray for ourselves when our soul has been struck by some grievous distress. Then we pray because we want something, because we desire something from God.

And if we were clear-sighted enough to see what is at stake—the dreadful ravaging of sin in the world, the threatening advance of the enemy—then intercession for our meetings, which we expect to be a battlefield where the powers of darkness and light meet in combat, would also become far more a cry of distress to the Lord and would bring victory from him.

Yet another matter may also be pointed out in this connection: united prayer, congregational prayer. It is scarcely ever practiced among us. Yes, someone says, when we are gathered for a prayer meeting, several of us pray for blessing upon the approaching edification meeting or mission meeting.

Now, yes, that is well and good where it is done. For not even this is practiced everywhere. But there may nevertheless be little true united prayer in it. It may be such, but it may also be only individual prayer, the prayer of single persons, and nothing more. For in the prayer for “blessing upon the meeting,” one person may have one thing in mind and another something entirely different. By united prayer for a meeting I mean particularly this: that a congregation has a living recognition of what it especially needs to receive from the Lord, and that it therefore, as one heart and one soul, longs to obtain this and therefore asks it of him in prayer.

“Again I say to you: Whatever two of you on earth agree to ask for, it shall be given them by my Father in heaven. For where two or three are gathered in my name, there I am in the midst of them.”

And concerning the apostles before the day of Pentecost it is said:

“They went up into the upper room, where they were staying—Peter and John and James and Andrew, Philip and Thomas, Bartholomew and Matthew, James the son of Alphaeus, and Simon the Zealot, and Judas the son of James. All these continued with one accord in prayer, together with some women and Mary, the mother of Jesus, and his brothers.”

Here we have true united prayer, the prayer that brings blessing, help from God. And for my part I am certain that if more prayer of this kind were practiced before our meetings, these would prove a greater blessing than they now often are.

Could we only remember vividly that the same Lord who has promised to hear prayer when two agree to pray for victory from God over evil—that he stretches out his nail-pierced hands toward human beings and says, “When I am lifted up from the earth, I will draw all to myself”—there would surely be more seriousness in our intercession. The conclusion, then, at which I arrive in the question of preparing our meetings by prayer is this: either no prayer at all is offered, or, if it is offered by one person or another, it is done with so little definite purpose that an answer from God to the prayer can scarcely be expected.

Certainly it is true that “it is God who works in us both to will and to accomplish,” so that our most heartfelt prayer, and that wrestling of the soul with him in which it is said, “I will not let you go until you bless me,” are also his work. Yet, in spite of this, he wills that we should pray with such earnestness as though the blessing of a meeting depended exclusively upon our prayer. But why can prayer be so significant when it too is his own impulse and work in us, brought forth by the Spirit?

My answer to this question is that our prayer means nothing other than that we are willing to receive the help and the blessing which the Lord, in free, unmerited grace, so gladly desires to give us, and which we so sorely need. If our soul has not been so deeply struck by our own and others’ spiritual distress that we ask him with full earnestness for deliverance in that distress, then this is the clearest sign that at bottom we see no danger and therefore do not care about deliverance from God in the distress. Then we shall never receive it. For salvation from sin and spiritual powerlessness is so great that it can be received only by a heart so pressed by this distress that, in comparison with it, it desires nothing else:

“Whom have I in heaven? And when I have you, I desire nothing on earth!”

In this cry we encounter the gathered mind, the undivided heart, for which salvation has become an overwhelmingly great reality, and which therefore partakes of it in so rich a measure.

What I intended by the little I have said here was, then, this: to point out a deficiency which, in my judgment, is so common at our larger meetings, namely, that as a rule they are far too little prepared for by united prayer, congregational prayer. And could I bring individual believers and our congregations to begin thinking over and considering this matter and, if possible, to take it more seriously, I should be deeply glad.

Another very common deficiency at most of our larger meetings is that those who are to take part by giving testimony often speak vaguely, unclearly, and without any aim.

And one of the most common causes of this is that the person concerned is unprepared. That is to say: the person has not studied the Word of God or the subject that is to be treated during the meeting.

But the natural consequence is that in such a case the speaker either does not speak on the text at all, or, if he keeps within the framework of the text, speaks so generally, so vaguely, and with so little point that the hearers are simply bored. And this ought not to be so when the word that is read is “spirit and life.” The Bible says: “Whoever speaks should speak as one who speaks words of God.” And God’s Word concerning the way of salvation is not religious phrases, but clear, definite, concrete truths about sin and guilt before God and about salvation from sin through Jesus Christ. And God’s Word also contains clear, definite truths concerning the life in God which his child is to live here on earth.

Therefore the testimony borne by layman or pastor during a meeting should consist of definite truths drawn from the subject or passage of Scripture that is being treated. No one has the right to excuse himself here with the often-repeated claim that one “has not had time to prepare.” For the one who has had no time for any preparation should not take part by giving testimony. If the matter concerning which we gather for a meeting is not so important that it drives us to as good and thorough a preparation as possible, it is better not to take up any of the meeting’s brief, precious time as a speaker. But not only the one who is to speak at a meeting should prepare himself. The one who is to hear should do the same: prepare himself through prayer and through consideration of the Word of God that is to be treated. By this means the benefit to the hearers would in every case be greater than it otherwise is.

The one who is to speak at a meeting should at least have some grasp of what the word on which he is to speak contains, and he should be able to set forth with reasonable clarity truths contained in the subject before him. Those who come to the meeting to hear have a right to expect this, and many of them really do expect it.

The time should at least be past when, at meetings such as those intended here, one can serve up old, familiar, assorted religious phrases without coherence and without any connection whatever with the subject of the meeting. Certainly this can still be done, and unfortunately it is probably done far too often; but the effects, the fruit, correspond to it. It is therefore careless and irresponsible when someone travels to a meeting to take part by giving testimony without having made any attempt to prepare himself. What would one say of the soldier who went into battle without having inspected his weapons and without having practiced using them? He would at least be able to accomplish little for the king or general in whose army he was to fight against the enemy.

And to take part by giving testimony at a meeting is to enter a holy war. It is to enter war against the prince of darkness and “the spiritual host of wickedness under heaven.”

Paul presents every Christian’s life here on earth as military service. He says:

“Stand, then, with truth girded about your loins and clothed with the breastplate of righteousness, and with your feet shod with the readiness that the gospel of salvation gives; and besides all this, take up the shield of faith, with which you shall be able to quench all the burning arrows of the evil one, and take the helmet of salvation and the sword of the Spirit, which is God’s Word.”

And then people will accuse Paul of not understanding what Christianity is.

But if this applies to Christians in general, then it applies to preachers in particular. For such men are not merely rank-and-file soldiers; they should be leaders, officers in Christ’s victorious army on earth. And where the armor and weapons are lacking or are not used, we need not wonder at the unsuccessful result of the battle.

None of us, therefore, should go to a meeting with the thought of accomplishing something for our Lord’s great cause without serious preparation through heartfelt prayer and a conscientious effort to penetrate the Word of God or the subject that is to be treated during the meeting.
